In this finite-horizon modular study, increasing width can enable delayed generalization and markedly accelerate the slow process after training interpolation is already possible. This rejects our strong width-independent decay-clock conjecture. The width relation changes with regularization and the task controls, and initial expected norm alone is insufficient wherever the matched-norm contrast persists. The supported explanation is conditional access to structured solutions under the complete training setup, with the particular circuit mechanism and eventual outcomes beyond the horizon still unresolved.
